% Draft manuscript insert. Confirm wording against the final, validated experiments.
\begin{abstract}
Compiler optimization quality can depend on the order in which optimization passes are applied, making pass ordering a search problem in its own right. This work presents Compiler Studio, an experimental framework that applies a Genetic Algorithm (GA) to search for compiler optimization-pass sequences over a simplified three-address-code intermediate representation. Six transformations are implemented: constant propagation, constant folding, copy propagation, common subexpression elimination, algebraic simplification, and dead code elimination. Candidate sequences are evaluated using a static instruction-cost model, and a Streamlit interface supports interactive experiments and inspection of optimized representations. The implementation includes both a fixed-length permutation baseline and an extended variable-length representation that permits passes to be omitted or repeated, fixed-point pipeline evaluation, input-driven straight-line examples, finite test-case output validation during candidate evaluation, multi-seed comparisons, pass ablation, and cost-weight sensitivity artifacts. Exhaustive enumeration remains feasible for the six-pass, one-use permutation space; therefore, those experiments are treated as a baseline rather than evidence that a GA is necessary. The current intermediate representation does not yet support branches or loops, and the static cost model is not a measurement of execution time. The experimental artifacts are intended to make these limitations explicit and to support reproducible follow-up evaluation.
\end{abstract}

\begin{IEEEkeywords}
Compiler optimization, phase ordering, genetic algorithms, evolutionary computation, intermediate representation, three-address code, program analysis, static cost model.
\end{IEEEkeywords}
